%  ____                 _ _
% |  _ \ ___  ___ _   _| | |_
% | |_) / _ \/ __| | | | | __|
% |  _ <  __/\__ \ |_| | | |_
% |_| \_\___||___/\__,_|_|\__|
\newpage
\section{Subject Morphing}
\label{sec:morphing}




% TODO: flip the sentences
Subject morphing involves blending semantics for generating images and videos. To our best knowledge, \textit{TrailBlazer} is  first in demonstrating subject morphing by prompt manipulation in the video diffusion domain. Related concepts have earlier been shown for image generation in MagicMix \cite{Magicmix} with, for example, the "corgi coffee machine".

While subject morphing through prompt embedding interpolation may seem less intuitive for real-world applications, it is widely used in the entertainment industry for example for superheroes (e.g., the She-Hulk can transform from a human to a monstrous character). For general usage, it could potentially serve as an entry point for generating new content that is more efficient than using a single prompt, particularly due to the limitations of CLIP \cite{CLIP}. For example, it might be challenging to generate a ``fish-like'' cat using the prompt "A fish-like cat walking on the grass" with a diffusion model. Instead, it would be easier to accomplish this goal by combining the prompt embeddings from ``A fish swimming on the grass'' and ``A cat walking on the grass.''

\FigRef{fig:morphing} illustrates the morphing outcomes generated by \textit{TrailBlazer}. All results are generated using default hyperparameter settings and involve linear interpolation of the prompt embeddings across video frames. The animated bounding boxes shift from right to left in the top two rows, and from left to right in the bottom two rows.

The outcome depicted in \FigRef{fig:morphing} demonstrates that morphing in a video clip can transition smoothly from one identity to another without significant artifacts. Notably, it avoids unrealistic deformations such as generating new joints in unexpected body parts (e.g., a tail on the head) or transforming one animal feature into another (e.g., an eye into an ear). Additionally, the subjects follow exactly the same motion in the synthesis (e.g., walking) across video frames.

\begin{figure*}[!htb]
  \centering
  \includegraphics[width=1.0\textwidth]{assets/supp-morphing.pdf}
  \caption{\textbf{Subject Morphing.} The prompts used starting from the first row: ``A [cat $\rightarrow$ dog] walking on the grass'', ``A [cat walking $\rightarrow$ fish swimming] on the grass'', ``A [parrot $\rightarrow$ king penguin] walking on the beach'', and ``A [tiger $\rightarrow$ elephant] walking in the wild park.''. Please refer to the text for more detail.}
  \label{fig:morphing}
\end{figure*}








\newpage
\section{Comprehensive ablations}
\label{sec:abla}

Given the limited space in the primary text, here we offer more supplementary ablation tests to substantiate our proposed approach. Broadly, we illustrate the impact of the spatial and temporal placement of guidance bounding boxes \emph{(bboxes)} on the overall result quality, exploring the effect of various bbox speed and size choices directed by user keyframing. To see details, please zoom in to the experiment images, and \textbf{especially refer to our supplementary video.}


\FigRef{fig:nodd} illustrates video synthesis using the pre-trained ZeroScope model
\textbf{without} applying our approach. Broadly, all the synthesized results exhibit fine details with plausible temporal coherence as would be seen in a real video featuring relatively slow motion. However, several side effects may be introduced alongside this realism. For example, the synthesized subject is often positioned in the same general area near the center of the images regardless of portrayed motion, and subjects like a galloping horse do not conveying the notion of speed. Additionally, artifacts such as extra or missing limbs (e.g., the cat in the second row) or other implausible results occasionally occur.


\begin{figure*}[!htb]
  \centering
  \includegraphics[width=1.0\textwidth]{assets/supp-no-dd.pdf}
  \caption{\textbf{Baseline results.} Each row shows equally-spaced frames sampled from a video generated using ZeroScope \emph{without applying our trajectory control approach}. The prompts used starting from the first row: ``A fish swimming in the sea'', ``The cat running on the grass field'', ``The horse galloping on the road'', and ``An astronaut walking on the moon''. These prompts are reused in subsequent examples in these supplementary results.}
  \label{fig:nodd}
\end{figure*}


\newpage
\subsection{Exploration and Ablation: Varied static bbox sizes}

\FigRef{fig:supp_static_bbox} shows the effect of the size of the bbox without considering motion.  The results indicate that the bbox size significantly influences the outcome. In extreme cases, the top row illustrates that a smaller bbox may yield unexpected entities in the area (e.g., white smoke next to the horse) or information leakage to the neighboring area (e.g., the blue attribute affecting the road). In contrast, the bottom row demonstrates that a overly large bbox can lead to broken results in general (e.g., the fish disappearing into the coral reef, and the strange blue pattern in place of the expected blue car). We expect this may be in large part due to the centered-object bias \cite{szabo-centeredobjects-21} in the pre-trained model's training data.

Our recommended bbox size falls within the range of 30\% to 60\% for optimal reconstruction quality. Note that very small- or large-sized bboxes can still be employed in our approach, but they are best specified for a particular frame rather than the entire sequence. This is demonstrated, for example, in \FigRef{fig:supp_dyn_bbox} guiding the swimming fish.

\begin{figure*}[!htb]
  \centering
  \includegraphics[width=1.0\textwidth]{assets/supp-static-bbox.pdf}
  \caption{\textbf{Static bbox sizes.} Each row shows the result of a static square bbox positioned at the center, where the width and height are 25\%, 50\%, and 90\% of the original image size (represented by the the green square on the left). The prompts used in the three sets of the experiments are: ``The \textbf{white horse} standing on the street'', ``The \textbf{fish} swimming in the sea'', and ``The \textbf{blue car} running on the road''.}
  \label{fig:supp_static_bbox}
\end{figure*}

\newpage
\subsection{Exploration and Ablation: Varied dynamic bbox sizes}

\FigRef{fig:supp_dyn_bbox} demonstrates video synthesis with a dynamically changing bbox size. In the top-left example, the bbox grows larger and then shrinks, resulting in a perspective effect where the fish swims towards the camera and then away from it. The frame highlighted in red indicates the middle keyframe with a large bbox. This aligns with our main text results in \FigRef{fig:main_dyn}, showcasing that the animated tiger and car respect the bbox size. The top-right example is a comparison to the top-left, portraying the fish only swimming toward the camera.

The second and the third rows show a comparison of the same bbox condition with the prompt words ``fish'' (second row), and ``sardine'' (third row), respectively. %Since ``fish'' and ``sardine'' represent different concepts in English, with ``fish'' being the super-set of ``sardine'',
This experiment aims to assess how well our method adapts to large bbox size variations, represented by the short/wide target bbox on the left and tall/thin target bbox on the right. The result on the left indicates that the output from the ``fish'' prompt does not adequately conform to the short-wide aspect ratio of the bounding box, whereas the result from the ``sardine'' prompt can more closely adjust to the desired bbox thanks to the elongated shape of the sardine.
Conversely, in the experiment on the right, both ``fish'' and ``sardine'' perform well with the tall/thin bounding box, since the tall aspect ratio can be satisfied by a fish facing directly toward or away from the camera. In general we expect that the obtained results will mimic the situations found in ZeroScope's training data,
while views that are outside the typical data (such as a fish swimming vertically, or a horse at the top of the image) will be difficult to synthesize.
%suggesting that the fish motion learned from ZeroScope generally involves normal swimming without ill-posed scenarios (e.g., backward swimming or lie-down-like swimming).

As with all our results, we see that the guided subject \emph{approximately} follows the specified bounding box, but does not exactly lie within the bbox.
While this is a disadvantage for some purposes, we argue that it is also an advantage for casual users -- if the subject exactly fit the bounding box it would require the user to imagine the correct aspect ratio of the subject under perspective (a difficult task for a non-artists) as well as do per-frame animation of the bbox to produce the oscillating motion of the swimming fish seen here.


\begin{figure*}[!htb]
  \centering
  \includegraphics[width=1.0\textwidth]{assets/supp-dyn-bbox.pdf}
  \caption{\textbf{Dynamic bbox sizes.} The result showcases six synthesized video sequences with the subject directed by the yellow arrow starting at the position indicated by green bbox. The number of the bboxes corresponding to the number of keyframes used in the experiment is, clockwise from top-left, $|\KKK| =$ 3, 2, 2, 2, 2, and 2, respectively. The prompt used in each result: ``The [X] swimming in the sea'', where ``[X]'' denotes the ``fish'' for the first and second rows, and ``sardine'' for the third row.}
  \label{fig:supp_dyn_bbox}
\end{figure*}

\newpage
\subsection{Exploration and Ablation: Speed control with multiple keys}

\FigRef{fig:supp_mul_keys} demonstrates controlling the subject's speed through varying the number of keyframes in the video synthesis. Given the recommended sequence length $N_f = 24$ for ZeroScope, we show the result of adding different
keyframes in between the start and end keyframes at the left/right image boundary, simulating the cat running back and forth on the grass field. It is clear that the cat moves relatively naturally according to the motion flow indicated by the yellow arrows. For instance, the cat looks back first before turning around, rather than showing an unnatural motion where the position of the head and tail is instantaneously swapped. As the cat moves faster, motion blur also introduced in the result annotated with red arrows. We found that this motion blur is hard to eliminate using negative prompts.

\begin{figure*}[!htb]
  \centering
  \includegraphics[width=1.0\textwidth]{assets/supp-mul-keys.pdf}
  \caption{\textbf{Speed Test: number of keyframes.} This result shows four synthesized video sequences with the cat's motion directed according to the yellow arrows starting from the position indicated by green bbox. The number of the arrows denotes the number of keyframes (excluding the start/end keyframes) used in each experiment. Specifically, starting from the top-left and proceeding in left/right top/down (English reading) order, there are
  % starting from top-left clockwise:
  $|\KKK| = $ 2, 3, 4, and 5, keyframes, respectively. The frames highlighted with red correspond to the user-specified keyframes,  excluding the start and end keyframes. The prompt used for all experiments is ``A \textbf{cat} running on the grass field''. The red arrows in the bottom-right example shows the introduced motion blur representing fast-moving speed.}
  \label{fig:supp_mul_keys}
\end{figure*}


\newpage
\subsection{Exploration and Ablation: Controlling speed with different placement of a single keyframe}

%Different from the results in \FigRef{fig:supp_mul_keys} at the previous page, \FigRef{fig:supp_speed_key} shows the speed control through the middle keyframe $\KKK_{f_1}$, where the first keyframe and the middle keyframe are with the same bbox value (e.g., $\BBB_{f_0} \equiv \BBB_{f_1}$) by adjusting its timing as indicated in the line chart at the second column.
\FigRef{fig:supp_speed_key} shows the results of moving the subject with increasing speeds. The first row shows the astronaut moving with constant speed obtained by the linearly interpolating bboxes at the left and right of the image.  Starting from the second row, the astronaut holds the position of the first bbox on the left side of the image for some period of time, then moves more rapidly to the right side of the image, as illustrated in the second column of the figure. This is obtained by changing the timing of a single ``middle'' keyframe $\KKK_{f_1}$, where the first keyframe and the middle keyframe have the same bbox location (e.g., $\BBB_{f_0} \equiv \BBB_{f_1}$).
 Similar to the results in \FigRef{fig:supp_mul_keys}, the synthesis may generate motion blur and artifacts when the speed is high (e.g., last row).

%In contrast to the outcomes in \FigRef{fig:supp_mul_keys} on the preceding page, \FigRef{fig:supp_speed_key} illustrates speed control through the middle keyframe $\KKK_{f_1}$, where the first keyframe and the middle keyframe share the same bounding box value (e.g., $\BBB_{f_0} \equiv \BBB_{f_1}$), with adjustments made to its timing as indicated in the line chart in the second column. The first row depicts the astronaut's linearly interpolated movement with a constant speed without the middle keyframe. From the second row onward, the astronaut maintains the position determined by $\KKK_{f_1}$, and then moves forward to the target bbox $\BBB_{f_2}$. Similar to the results in \FigRef{fig:supp_mul_keys}, higher speeds (e.g., last row) may lead to motion blur or undesirable entity generation in the output.

\begin{figure*}[!htb]
  \centering
  \includegraphics[width=1.0\textwidth]{assets/supp-speed-key.pdf}
  \caption{\textbf{Speed Test: the timing of a keyframe.} The result shows four synthesized video sequences with the subject directed according to the yellow arrow starting at the position indicated by green bbox, as illustrated in the first column. All experiments except the first use three keyframes ($|\KKK| =3$), where the timing of the internal keyframe (e.g., $\KKK_{f_1}$) controls the duration of a stationary phase and the speed of the subsequent motion,   as illustrated in the second column.
  The horizontal and vertical axis in the second column represent the left/right position and timing, respectively. The frame outlined in red indicates the frame controlled by $\KKK_{f_1}$, corresponding to the time when the astronaut starts to move. The prompt used for all experiments: ``The \textbf{astronaut} walking on the moon''.}
  \label{fig:supp_speed_key}
\end{figure*}

\newpage
\subsection{Exploration and Ablation: Irregular trajectory}


%We demonstrate the irregular trajectory determined by varied keyframe in \FigRef{fig:supp_irr_path}. The four experiments with the trajectory of zigzag (top-left), triangle (top-right), discontinuity (bottom-left), and the down-pointing triangle (bottom-right). All horse result showcases the fast speed running with motion blur. However, the results with turning points has limitation to show the horse turning around and creates the artifacts. E.g., the third frame of down-pointing triangle case, that the horse swaps its head and tail. We think it is somewhat expected because it is naturally that a horse is good at charging forward rather than turning around with high speed like the cat or dog. In summary, the consistency between the prompt, the timing and bboxes in the keyframe must be maintained to increase the result fidelity.

We illustrate irregular trajectories determined by varied keyframes in \FigRef{fig:supp_irr_path}. The four experiments involve a zigzag trajectory (top-left), a triangle trajectory (top-right), a \emph{discontinuous} trajectory (bottom-left), and a down-pointing triangle trajectory (bottom-right). In every result the horse shows high-speed running with motion blur. However, the results with turning points show limitations in depicting the horse quickly turning around and may show artifacts. For example, in the third frame of the down-pointing triangle case, the horse appears to swap its head and tail. Difficulty portraying this turn is somewhat expected, as horses cannot naturally execute tight high-speed turns, unlike cats or dogs. On the other hand, the down-pointing triangle video naturally introduces a perspective-like size change as the horse moves higher in the image,
%Notably, some results show the perspective (e.g., down-pointing triangle) that the horse charging to the camera and then moving away,
similar to the previous results in \FigRef{fig:supp_dyn_bbox}, and also the tiger example \FigRef{fig:main_dyn} in our main text. In summary, maintaining consistency between the prompt and the timing and location of the keyframed bounding boxes is crucial for producing realistic results.

\begin{figure*}[!htb]
  \centering
  \includegraphics[width=1.0\textwidth]{assets/supp-irr-path.pdf}
  \caption{\textbf{Irregular trajectory.} The figure shows four synthesized video sequences with the horse subject directed according to the yellow arrows starting from the position indicated by green bbox.
  % jp think the 2,3,4,5 here is wrong, maybe not needed since is explained in the previous multiple keys figure
  %The number of the arrows denotes the number of internal keyframes (excluded start/end keyframes) used in the experiment, starting from top-left clockwise: $|\KKK| =$ 2, 3, 4, and 5, respectively.
  The frames highlighted in red correspond to keyframes. The start and end keyframes are not indicated. The prompt used for all examples: ``A \textbf{horse} galloping on the road''.}
  \label{fig:supp_irr_path}
\end{figure*}
